\chapter{系统协同：目标、身份与执行控制的联合动力学}
\label{chp:system_synergy}
本章追问智能系统在不存在单一中央控制者时如何形成连续、可审计而又可修订的行动。上一章已经把体验描述为偏好、风险、身份历史与任务度量对比较和搜索的影响；本章进一步区分目标规范、身份状态和执行控制，研究它们怎样在不同时间尺度上联合，而不把三者等同为场、孤立子和热力学泵。目标说明何为可接受结果，身份说明谁承担承诺与后果，执行控制把候选变成受边界约束的动作；三者都依赖环境回执，也都可能失配。我们将以带类型状态、版本、权限和资源账本的离散动力学重建原有三体图景，并明确内部误差下降、动作提交和真实任务成功之间的边界。

\section{对象分工：目标规范、身份状态与执行控制}
\label{sec:section_9326}
历史文本把目的、自我和宏观层分别称为场、核与泵。这个分工抓住了方向、连续性和行动三个功能，却把不同类型的对象放进同一物理图景。我们改用三类计算对象。目标规范记为 $\Theta_k=(\mathcal C_k,J_k,\mathcal R_k,\nu_k)$：$\mathcal C_k$ 是不可违反的可行域，$J_k$ 是候选轨迹的任务代价，$\mathcal R_k$ 是风险泛函，$\nu_k$ 是版本与来源。身份状态记为 $S_k=(I_k,H_k,K_k,P_k)$：身份键 $I_k$ 维持对象归属，历史 $H_k$ 保存带来源事件，承诺 $K_k$ 约束跨时行动，权限 $P_k$ 界定谁能改写哪些组件。执行控制记为 $U_k=(\pi_k,B_k,A_k,R^u_k)$，分别表示策略、边界检查器、动作接口和资源状态。三者空间、单位和更新频率不同，不能用乘法口号替代接口定义。

目标规范并不直接“弯曲”所有表示。它通过任务相关读出选择状态分量、比较候选并约束提交。一个学习目标可以改变排序，却不能凭自身创造环境中不存在的动作；一条硬边界可以淘汰动作，却不负责给剩余候选排序。目标还可能来自系统稳态、用户指令、组织制度或环境约束，来源冲突需由显式优先序和授权解决。若只保留一个奖励标量，用户临时偏好可能覆盖安全边界，环境回执也可能被语言上的自我确认取代。我们因此分开记录代价、风险、约束和权限，并要求每次目标切换携带版本。

身份状态提供跨任务连续性和归责，而不是所有运动的绝对坐标原点。系统可以在没有丰富自我模型时完成局部反射或一次性工具调用；只有长期承诺、个人化历史、角色切换和责任追踪需要较强身份接口。身份也不是不可改变的拓扑不变量。新证据、授权和能力变化可以产生候选版本，但必须通过归属核验、承诺回归、安全检查和回滚准备后提交。把一切抗改写都解释为保护自我，会把错误信念和权限滥用包装成稳定性；把一切适应都视为身份更新，又会使系统随强提示漂移。

执行控制负责把状态估计和目标规范转化为动作包。它包含候选生成、评分、边界裁剪、资源调度、动作提交和环境确认，任何一步失败都不能被“宏观做功”一词掩盖。算法计算量、等待时间、设备焦耳消耗和任务代价分别记账。系统可能知道目标且保持身份，却因权限不足、工具故障或资源不足无法行动；也可能有强执行能力，却因目标版本错误而高效地做错事。执行资格因此是条件性的，并由外部回执而非内部激活强度最终确认。

三者通过类型化接口相连。任务读出 $R_{\Theta_k}(h_k,S_k)$ 生成带理由的候选，边界函数 $B_k$ 检查可行性，策略 $\pi_k$ 在资源预算内排序，提交器把动作发送给环境，结果 $y_{k+1}$ 再分别更新信念、目标证据和身份历史。局部结构网络保存身份、来源、承诺和权限，全局分布表示提供相似性和候选；它们通过绑定接口联合。AgentOS可用CCM协调目标、TCE计算控制、SPC保存结构、Boundary执行约束、$M_e$保存证据，但这些只是工程实例，不是所有智能必须采用的唯一机制。

验收先做分离干预：只改变目标版本，应改变相关候选排序而不改身份键；只切换角色，应改变权限与承诺读出而不重写环境事实；只限制资源，应影响响应时间和计划深度而不伪造价值变化；只断开动作接口，应保留意图但不声称任务完成。再做组合测试，检查接口延迟、冲突仲裁、来源伪造和回滚。只有这些可观察差异成立，我们才说目标、身份与执行控制形成了系统协同，而不是三个物理名称共享一幅示意图。

\section{耦合方程：从目标错配到受约束闭环控制}
\label{sec:section_4740}
原稿用“意志等于自我乘以目的减本能”概括冲突。计算上，我们先定义错配的对象。令 $h_k\in\mathbb R^{d_h}$ 为当前估计，$g_{\Theta_k}(h_k)\in\mathbb R^{d_g}$ 为任务读出，$r_k\in\mathbb R^{d_g}$ 为允许随版本变化的目标参考，$W_k\succeq0$ 为无量纲或已校准的比较矩阵。即时错配为 $e_k=g_{\Theta_k}(h_k)-r_k$，但它只描述所选读出上的差异，不等于痛苦、能量或失败。身份状态通过权限、承诺与风险权重影响哪些误差可被控制，而不是作为内积中的神秘放大因子。

在预测窗口 $H$ 内，我们把控制问题写为
\[
\min_{u_{k:k+H-1}}\sum_{j=0}^{H-1}
\bigl(e_{k+j}^{\top}W_{k+j}e_{k+j}+c_u(u_{k+j})+\lambda R_{k+j}\bigr),
\]
并要求状态转移 $h_{j+1}=F(h_j,u_j,o_j;G_j)$、动作权限 $u_j\in\mathcal U(S_j,P_j)$ 和边界约束 $b_m(h_j,u_j)\leq0$ 同时成立。这是一个模型预测控制式工程设计，不是普遍意志定律。目标、结构、输入和度量变化时，$r_k$、$G_k$、$o_k$ 与 $W_k$ 必须进入更新；离散步长、窗口和求解容差也要报告。若问题不可行，控制器应返回冲突证明或降级方案，而不是注入更大“力量”。

想玩游戏但需要学习的例子可以据此拆解。游戏候选可能具有较低即时努力和较高短期偏好，学习候选可能更符合考试目标和长期承诺。系统先核验考试期限、休息需求、用户授权和当前疲劳，再比较两类轨迹。若连续学习已越过健康边界，继续学习并非更有意志；若考试目标已取消，旧承诺也不应永久占优。可靠控制可以安排二十五分钟学习、五分钟休息，再用完成度和状态回执滚动重规划。这里的“意志输出”是动作序列、边界理由和更新日志，不是未定义的矢量力。

目标下降还不能推出任务成功。内部读出可能漏掉真实约束，环境模型可能错误，工具可能执行失败，用户目标也可能含混。每次动作后都要等待环境确认 $y_{k+1}$，通过观测模型更新 $h_{k+1}$，并区分预测误差、执行误差和规范冲突。若模型连续预测“已学习”但环境中没有笔记、测试成绩或可复演成果，内部损失下降只是自洽。对于不可逆动作，回滚只能恢复内部版本，不能撤销已发生的环境后果，因此提交前需要更严格的模拟、授权和风险预算。

身份权重的作用也应受限。与长期承诺高度相关的任务可以获得更高优先级，但身份叙事不能越过安全与他者权利。我们令门控 $m(S_k,\Theta_k)\in[0,1]$ 调整软目标权重，并对变化率和最大增益设界；硬约束另行处理。这样可以避免“关乎尊严”被用来无限放大动作。多主体系统还需标明提议者、批准者、执行者和后果承担者，不能用一个群体自我掩盖责任。

稳定性只在声明条件下讨论。若固定结构和目标窗口内 $F$ 对状态局部Lipschitz，控制增益有界，延迟小于给定阈值，并且离散步长满足数值条件，我们可以分析局部有界或输入到状态稳定；目标频繁切换、结构事件、不可观测状态和求解超时会使结论失效。工程测试因此包括阶跃目标、冲突目标、延迟回执、动作故障、预算耗尽和恶意身份提示。验收同时报告误差轨迹、约束违例、资源、动作成功率和环境结果，不能用一个“意志强度”总分替代。

在这一重建中，目的提供版本化规范，身份提供连续性与权限，执行控制提供受约束动作。三者的联合来自消息、门控、验证和回执，而非严格的规范场耦合。下一节将沿用这一分工，重新分析原稿所谓机会主义、虚无和瘫痪，把它们限定为可诊断的工程失配，而不借系统故障对人类临床状态作物理化归类。
